% Reproducibility checklist: the paper's supplement, outside its pages, built by
% paper/supplement.tex (which defines \file). Results are macros from generated/numbers.tex;
% settings, versions and commands are those of the files named in the last column and of
% docs/RUNBOOK.md. The repository is private, so nothing here links to it.
\section*{Reproducibility checklist}

The code, the settings, the result summaries, the frozen pre-registration and the deviation
log are kept in one private repository, available from the authors on request; the EEG data
are public. Paths below are relative to that repository, and commands run from its top
folder; \file{docs/RUNBOOK.md} says what each one should print. Classical methods give
identical numbers on a re-run; networks trained on a GPU can differ slightly.

\begin{center}
\footnotesize
\begin{tabularx}{\textwidth}{@{}>{\raggedright\arraybackslash}p{2.3cm}>{\raggedright\arraybackslash}X>{\raggedright\arraybackslash}p{4.4cm}@{}}
\toprule
Item & What was used & Where it is stated \\
\midrule
Data and version & The authors' preprocessed EEG of \nSubjects{} stroke patients
(\file{processeddata.zip}, \nRuns{} runs): Wang et al., \emph{Scientific Data} 13, 1109
(2026), doi:10.1038/s41597-026-07787-y; figshare doi:10.6084/m9.figshare.30509165,
version~11, CC~BY~4.0; checked against the published MD5 and every file's size. H1b:
BNCI2014\_001 and Cho2017, downloaded by MOABB 1.7.2 &
\file{configs/data.yaml}, \file{docs/dataset-card.md}, \file{results/audit/inventory.csv} \\
Integrity audit & Every run checked for format, flat channels and recordings that appear
twice; \nRunsKept{} runs in \nSessionsKept{} sessions kept, sub-21 and sub-22 linked &
\file{results/audit/audit_report.md}, \file{configs/data.yaml} (exclusions) \\
Preprocessing & The authors' band-pass, reference and artefact rejection as published;
40 of 43 channels (Fp1, Fp2 and Fpz dropped); a window from 1 to 5\,s after every rest
marker and cue, at \nSamplingRate\,Hz; windows that cross a removed segment or a run's end
dropped (\nWindowsDropped), \nWindows{} kept; cached with a manifest that every later step
checks & \file{configs/data.yaml}, \file{results/audit/epoch_counts.csv},
\file{results/audit/epoch_drops.csv} \\
Splits & Leakage-free cross-validation: 10 folds $\times$ 10 repeats within a session,
every step fitted inside the training folds. All else: learn on runs 1--2, test on runs 3--4, which no step sees.
Calibration draws $k \in \{0, 1, 2, 3, 5, 10, \text{all}\}$ trials per class from runs 1--2,
ten random draws (five for networks) shared by every method, plus one chronological draw.
Sources: every other patient, never the target's linked partner. A check before every fit
stops the run if a test trial enters training or the target or its partner enters the
source data & \file{configs/protocol.yaml}, \file{src/sixthfinger/evaluation/},
\file{tests/test_leakage.py} \\
Hyperparameters & Fixed before the results, one file per method. The three open transfer
values were chosen once by a pre-registered rule on development patients sub-01 and sub-03,
learning on run 1 and scoring on run 2, never on their runs 3--4; the sources were, as in the
main study, every other patient, test runs included. Chosen: a target weight of
\rqTwoHparamRctTslrWeight{} of \{1, 5, 20\} (development means
\rqTwoHparamRctTslrMeanMin--\rqTwoHparamRctTslrMeanMax\,\%, three classes); EEGNet fine-tuning
at \rqTwoHparamEegnetLr{} for \rqTwoHparamEegnetEpochs{} epochs & \file{configs/models/*.yaml},
\file{scripts/select_transfer_hparams.py} \\
Seeds & Folds and calibration draws from seed 0 (draw $i$ uses seed $i$); bootstrap
10\,000 resamples with seed 12345; EEGNet from scratch with seeds 0, 1 and 2, their mean
reported; pre-training with seed 0 & \file{configs/protocol.yaml},
\file{configs/experiments/dl_baselines.yaml}, \file{configs/models/eegnet_ea_ft.yaml} \\
Software & Python 3.12; NumPy 2.5.3, SciPy 1.18.1, pandas 3.0.6, MNE 1.13.2, scikit-learn
1.9.1, pyRiemann 0.12, PyTorch 2.11.0, braindecode 1.8.1, ONNX 1.23.1, onnxruntime 1.30.0,
pingouin 0.7.0, statsmodels 0.15.0, MOABB 1.7.2; the PDFs with Tectonic &
\file{pyproject.toml} and \file{uv.lock} (every version pinned) \\
Hardware & An \demoInferenceMachine{} laptop with 8\,GB of memory: classical methods,
follow-ups, the demonstrator's timing. Kaggle notebooks: a \dlPretrainDevice{} GPU for every
EEGNet of the main study, CPUs for some classical calibration shards &
\file{results/summary/dl_summary.json}, \file{results/summary/demo_inference_timing.json} \\
Run times & One CSP+LDA fit in leakage-free cross-validation: \runtimeCspLdaCvLeakfreeFitSeconds\,s
(\runtimeCspLdaCvLeakfreeFits{} fits). EEGNet pre-training: \dlPretrainMinutes{} GPU minutes
for \dlPretrainModels{} models (\dlPretrainJobMinutes\,min wall clock). Healthy-grasp
pre-training (exploratory): \exThreeArmGpPretrainHours\,h on the laptop's GPU. One
demonstrator decision: \demoInferenceMedianMs\,ms. Other run times were not recorded as
results; the runbook gives estimates & \file{results/summary/rq1_runtime.csv},
\file{dl_summary.json}, \file{ex3_grasp_pretraining_checks.json},
\file{demo_inference_timing.json}, \file{docs/RUNBOOK.md} \S4 \\
Pre-registration & Hypotheses, protocol and analysis frozen in git on a commit of
2026-10-03, tagged \file{prereg-v1} on 2026-10-08, before any confirmatory result; results
frozen under \file{results-v1} (2026-10-12); every later change is one of \nDeviations{}
dated entries of the deviation log, and analyses added after the results are exploratory.
The dates are git's own: the plan was not registered with an independent registry &
\file{docs/preregistration/} (a verbatim copy of both documents, with their provenance), git
tags \file{prereg-v1}, \file{results-v1} \\
Provenance & Every result row stores the commit that made it and hashes of its settings and
of the window manifest; every number in the paper is generated from the result summaries &
\file{scripts/make_paper_numbers.py}, \file{paper/generated/numbers-index.md} \\
Code & MIT licence; available from the authors on request & \file{LICENSE},
\file{CITATION.cff} \\
\bottomrule
\end{tabularx}
\end{center}

\subsection*{Deviation log}

Every change after the pre-registration's tag is an entry of the deviation log, with its date,
its reason and when it was made relative to the results. The \nDeviations{} entries are listed
below, one line each in our own words; the full text of each is in
\file{docs/preregistration/deviations.md}, available from the authors on request.

\begin{center}
\footnotesize
\begin{tabularx}{\textwidth}{@{}p{1.6cm}>{\raggedright\arraybackslash}X>{\raggedright\arraybackslash}p{4.6cm}@{}}
\toprule
Date & Entry & When \\
\midrule
2026-10-09 & RQ1b to run on BNCI2014\_001 alone, for lack of disk space (superseded the same day) & before any public data or RQ1b result \\
2026-10-09 & Input of the rebuilt authors' recipe: windows whose wider span crosses a cut are dropped & before any RQ1 result \\
2026-10-09 & FBCSP's feature selection gets a fixed random seed (a clarification) & before any RQ1 result \\
2026-10-09 & Cho2017 run as pre-registered; the row above is superseded & before any Cho2017 data \\
2026-10-10 & Rebuilt recipe: a fold where unregularised CSP has no solution is recorded as missing & during RQ1; H1 unaffected \\
2026-10-11 & $k = \text{all}$ takes every trial of the pool & before any calibration result \\
2026-10-11 & The fine-tuned EEGNet keeps its three-output network for every task & before any calibration result \\
2026-10-11 & Choosing the transfer values: what ``no runs 3--4 are touched'' means, and the tie order & before any calibration result \\
2026-10-11 & The target's reference for re-centring and alignment is the draw's unlabelled set & before any calibration result \\
2026-10-11 & Re-centred methods computed with our own re-centring (same operations as pyRiemann) & before any calibration result \\
2026-10-11 & Scope of the post-therapy replication and of the chronological draw & before any calibration result \\
2026-10-11 & Secondary EEGNet pair: the tests of H2a, H2b and H3 on the pre-therapy session & before any calibration result \\
2026-10-11 & $k^*$, trials saved and their units & before any calibration result \\
2026-10-11 & Descriptive summaries outside the pairs & before any calibration result \\
2026-10-11 & The no-signal check with permuted labels & before any calibration result \\
2026-10-11 & When the compute contingency for EEGNet triggers & before any calibration result \\
2026-10-11 & A sanity check of the calibration study against the benchmark & before any calibration result \\
2026-10-11 & Disclosure of the results already known before the calibration study & before any calibration result \\
2026-10-11 & ``Exact $p$'' means the numeric $p$-value & after RQ1 results; no category changes \\
2026-10-11 & The log is also copied into the pre-registration (bookkeeping) & after the tag \\
2026-10-12 & Draws below all pick only trials that kept both windows & before any calibration result \\
2026-10-12 & EEGNet baselines and pre-training re-run: the last incomplete batch had been skipped & before any calibration result \\
2026-10-12 & The configuration hash ignores settings left at their default & before any calibration result \\
2026-10-12 & Compute contingency not triggered & before the full EEGNet runs \\
2026-10-12 & The sanity check of 2026-10-11 failed as defined; its premise was wrong & after the calibration results; no confirmatory value changes \\
2026-10-12 & Exploratory study reduced: no electrode-count study; mirroring with classical methods only & after the calibration results, before any exploratory result \\
2026-10-12 & Two descriptive analyses added, labelled post hoc: the channel-group check and replay details & after the calibration results \\
2026-10-12 & The demonstrator's patient: breaking a tie in the rule & before the demonstrator's bundle \\
2026-10-23 & Exploratory: more data per decision for hand against sixth finger & after every confirmatory result \\
2026-10-23 & Exploratory: the raw recordings for hand against sixth finger & after every confirmatory result \\
2026-10-23 & Exploratory: pre-training with healthy grasp imagery & after every confirmatory result \\
2026-10-23 & Exploratory: the sixth finger as its own command (the switch) & after every confirmatory result \\
2026-10-23 & Exploratory: the switch played as if live & after every confirmatory result \\
\bottomrule
\end{tabularx}
\end{center}

\subsection*{Commands}

\begin{center}
\footnotesize
\begin{tabularx}{\textwidth}{@{}>{\raggedright\arraybackslash}p{3.0cm}>{\raggedright\arraybackslash}X@{}}
\toprule
Step & Command \\
\midrule
Environment & \texttt{make setup}; check with \texttt{uv run sixthfinger doctor} \\
Data, audit, windows & \texttt{make data}; \texttt{make verify-data}; \texttt{make audit};
\texttt{make epochs} \\
Tests & \texttt{make check} \\
Leakage (H1) & \texttt{uv run sixthfinger run configs/experiments/rq1\_benchmark.yaml}, the
same with \texttt{rq1\_context.yaml}; \texttt{uv run python scripts/summarise\_rq1.py} \\
Healthy datasets (H1b) & \texttt{uv run -{}-extra moabb python scripts/rq1b\_public.py
-{}-datasets BNCI2014\_001}, the same with \texttt{Cho2017}; then
\texttt{-{}-summarise-only} \\
EEGNet baselines and pre-training & \texttt{bash notebooks/kaggle/bundle.sh -{}-code};
\texttt{bash notebooks/kaggle/push.sh dl-baselines gpu run
configs/experiments/dl\_baselines.yaml}; \texttt{bash notebooks/kaggle/push.sh pretrain gpu
pretrain}; \texttt{bash notebooks/kaggle/fetch.sh <name>}; \texttt{uv run python
scripts/merge\_scores.py results/kaggle/*/results/runs}; \texttt{uv run python
scripts/summarise\_dl.py} \\
Transfer values & \texttt{uv run python scripts/select\_transfer\_hparams.py -{}-methods
rct\_tslr rct\_mdm eegnet\_ea\_ft} \\
Calibration (H2, H3) & \texttt{uv run sixthfinger run configs/experiments/<name>.yaml} for
\texttt{rq2\_calib\_s1}, \texttt{rq3\_calib\_s2}, \texttt{rq2\_rq3\_post} and
\texttt{rq2\_nosignal} (the EEGNet methods through \texttt{push.sh <name> gpu -{}-models run
...}); \texttt{uv run python scripts/summarise\_calibration.py};
\texttt{uv run python scripts/summarise\_calibration.py -{}-minutes-only} (the
recorded minutes) \\
Exploratory analyses & \texttt{rq4\_mirror.yaml} through \texttt{push.sh}, then
\texttt{scripts/summarise\_rq4.py}; \texttt{scripts/rq4\_longitudinal.py};
\texttt{scripts/pseudo\_online.py}, then \texttt{-{}-summarise};
\texttt{scripts/check\_visual\_cue.py}; \texttt{scripts/make\_explain\_figures.py}; the
follow-ups \texttt{scripts/ex1\_transfer\_ah\_sf.py} to \texttt{ex5\_live\_switch.py}, one
\texttt{-{}-part} at a time in the order each script's docstring lists \\
Demonstrator & \texttt{make demo-bundle}; \texttt{uv run python
scripts/time\_demo\_inference.py}; \texttt{uv run python scripts/count\_eegnet\_weights.py} \\
Numbers and PDFs & \texttt{uv run -{}-with
brotli==1.2.1 python scripts/make\_paper\_figures.py} and the same for
\texttt{scripts/make\_rq3\_patient\_figure.py}; \texttt{uv run python
scripts/make\_paper\_numbers.py}; \texttt{make paper}; this supplement with \texttt{make
paper-supplement} \\
Checking a re-run & \texttt{uv run python scripts/compare\_summaries.py <commit>}: exit
code 0 when no summary changed \\
\bottomrule
\end{tabularx}
\end{center}
